\uuid{AJln}
\titre{Fraction}
\auteur{bodin ; hanani ; mzari}
\organisation{exo7}
\licence{CC BY-NC-SA 4.0}
\datecreate{2020-01-20}
\niveau{L1}
\module{Analyse}
\chapitre{Propriétés de R}
\sousChapitre{Les rationnels}
\difficulte{2}
\theme{Rationnels, fraction}
\exo7qcmid{lille-1.reels.003}

\contenu{
\begin{question}
\qclassification{120.01}
\qtags{motcle=fraction}
Soient $x$ et $y$ deux nombres rationnels strictement positifs.
Parmi les nombres réels suivants, lesquels sont aussi des nombres rationnels ?
\begin{answers}
    \good{$\frac{x-y}{x+y}$}

    \bad{$\frac{\sqrt{x}}{\sqrt{y}}$}

    \good{$x-y^2$}

    \good{$(\sqrt{x}-\sqrt{y})(\sqrt{x}+\sqrt{y})$}

\end{answers}
\begin{explanations}
La somme, le produit, le quotient de deux nombres rationnels reste un nombre rationnel. La racine carrée d'un nombre rationnel n'est pas toujours un nombre rationnel (par exemple la racine carrée de $2$). Par contre, par identité remarquable, on a $(\sqrt{x}-\sqrt{y})(\sqrt{x}+\sqrt{y}) = x^2-y^2$ qui est un nombre rationnel.
\end{explanations}
\end{question}
}
